% Publication adaptation of I236, frozen SHA256 a30077ca36f40e6613e340d3f8ae0ac692000f7b39d6eb4113897f40b7c68931.
% New presentation requires its own review; this file is not the reviewed original.
\section{Finite-viscosity fitted-class rigidity}\label{app:viscous}

This appendix supplies the positive-viscosity input used in the main
argument. Its proof of record is I236, \emph{analytic-attempt},
Sections 1--6, previously reviewed in I239. The full argument is
carried here rather than replaced by that review verdict.
It includes the contact-propagation proof; it does not assume
regularity of arbitrary entire fitted solutions. Its geometric inputs
are the fit compactness and barycentric-envelope lemmas of
\cite[Sections 4--5]{OpenAI}, with the equation changes checked below.

Fix an integer $m\geq1$ and $0<\tau<\infty$. For $P=(p,q)\in
\mathbb R^m\times\mathbb R$ and a symmetric matrix $X$, define
\[
 A_\tau(p)=\begin{pmatrix}\tau I+pp^T&p\\p^T&1\end{pmatrix},
 \qquad G_\tau(P,X)=\operatorname{tr}(A_\tau(p)X).
\]
The time-gradient component $q$ does not enter $G_\tau$. We use the
positive-Hessian convention: upper tests for subsolutions give
$G_\tau\geq0$, and lower tests for supersolutions give $G_\tau\leq0$.

\begin{theorem}[Finite-viscosity rigidity]\label{thm:viscous}
Let $V\in C(\mathbb R^m\times\mathbb R)$ solve $G_\tau=0$ in the
viscosity sense. Suppose a finite $C_0\geq0$ has the following property:
for every $(x,t)$ and every $R>0$, some $p\in\mathbb R^m$ and
$a\in\mathbb R$ satisfy
\begin{equation}\label{eq:app-fit}
 \left|V(x+h,t+\ell)-V(x,t)-p\cdot h
       -\left(a-\frac{|p|^2}{2}\right)\ell\right|\leq C_0R
 \quad (|\ell|\leq R,\ |h-p\ell|\leq R).
\end{equation}
Then $V$ is affine. No gradient bound or energy assumption is required.
No estimate uniform as $\tau$ approaches zero or infinity is asserted.
\end{theorem}

\subsection{Time regularity of spatially semiconcave subsolutions}

\begin{lemma}\label{lem:time-lipschitz}
A continuous $G_\tau$-subsolution with locally uniform spatial
semiconcavity is locally Lipschitz in all variables.
\end{lemma}
\begin{proof}
Work on a product cylinder compactly contained in the domain and a
smaller cylinder of base points $(x_0,t_0)$. Choose $r>0$ so that every
$\overline B_r(x_0)$ stays inside the larger spatial ball. Enlarge the
semiconcavity constant to $K\geq0$. Continuity supplies a common bound
$M$ and a time modulus $\omega$ on the larger compact cylinder.

At every base point, a supporting hyperplane to
$u(\cdot,t_0)-K|\cdot|^2/2$ gives a vector $p$ with
\[
 u(y,t_0)\leq u(x_0,t_0)+p\cdot(y-x_0)+\frac K2|y-x_0|^2
 \quad (|y-x_0|\leq r).
\]
Evaluating at opposite fixed increments bounds $|p|$ by a common $L$.
The same argument at neighboring points bounds spatial Lipschitz
constants uniformly in time. Set $K'=K+1$ and hold this $p$ fixed in
\[
 h(t)=\max_{|y-x_0|\leq r}
 \left\{u(y,t)-p\cdot(y-x_0)-\frac{K'}2|y-x_0|^2\right\}.
\]
At $t_0$ the unique maximizer is $x_0$, and boundary competitors have
a deficit at least $r^2/2$. Choose a common $s>0$ with
$\omega(s)<r^2/8$ and with the relevant time intervals inside the
larger cylinder. For $|t-t_0|<s$, boundary competitors are at most
$u(x_0,t_0)-3r^2/8$, whereas the central competitor is at least
$u(x_0,t_0)-r^2/8$. Thus every maximizing point is interior. Also
$h$ is continuous and $|h|\leq H:=M+Lr+K'r^2/2$.

If a scalar upper test $\psi$ touches $h$ at $t$ and $y$ is a maximizer,
then
\[
 \phi(Y,T)=\psi(T)+p\cdot(Y-x_0)+\frac{K'}2|Y-x_0|^2
\]
is an upper test of $u$ at $(y,t)$, after adjusting the contact constant.
Its mixed Hessian is zero. The subsolution inequality gives
\[
 0\leq\psi''(t)+K'\bigl(|p+K'(y-x_0)|^2+\tau m\bigr),
 \qquad \psi''(t)\geq-C,
 \quad C=K'\bigl((L+K'r)^2+\tau m\bigr).
\]
Consequently $h(t)+C(t-t_0)^2/2$ is convex. The one-dimensional
viscosity characterization used here follows directly from the chord
criterion: failure of convexity yields an interior upper contact by a
strictly concave quadratic added to an affine function.
The slopes of this bounded convex function are bounded on the middle
half interval by its endpoint secants. For example, $8H/s+4Cs$ bounds
the Lipschitz constant of $h$ there. Since
$u(x_0,t)\leq h(t)$ with equality at $t_0$, this bounds
$u(x_0,t)-u(x_0,t_0)$ from above. Repeating at base time $t$, with the
same constants, gives the reverse bound. Combining with the spatial
bound proves the lemma.
\end{proof}

\subsection{Preserving contact while making the supersolution Lipschitz}

Suppose $u\leq v$, $u$ is locally Lipschitz, $v$ is a continuous
supersolution, and $u(z_0)=v(z_0)$. Work inside a ball with a compact
larger margin. If $L_0$ bounds the Lipschitz constant of $u$ there,
choose $L>L_0$ and a small $r>0$. Define locally
\[
 v_L(z)=\min_{|h|\leq r}\{v(z+h)+L|h|\}.
\]
The minimum is attained, and $v_L$ is continuous. Every competitor
obeys
\begin{equation}\label{eq:cone-lower}
 v(z+h)+L|h|\geq u(z)+(L-L_0)|h|.
\end{equation}
Thus $u\leq v_L\leq v$, and the unique minimizer at $z_0$ is zero.
For any minimizing $h$,
\[
 (L-L_0)|h|\leq v(z)-u(z).
\]
Continuity therefore places all minimizers in $B_{r/2}$ on a small
neighborhood of $z_0$.

A lower test $\phi$ of $v_L$ at $z$, pulled back by a minimizing
fixed translation $h$, is a lower test of $v$ at $z+h$, with unchanged
gradient and Hessian. The autonomous equation implies that $v_L$
remains a supersolution. If $h$ minimizes at $z$, then
$h+z-z'$ is admissible at a sufficiently close $z'$, so
\[
 v_L(z')\leq v(z+h)+L|h+z-z'|
           \leq v_L(z)+L|z-z'|.
\]
Reverse the roles of the two points to obtain local $L$-Lipschitz
continuity. Finally, if $v_L(z)=u(z)$, equation
\eqref{eq:cone-lower} forces its minimizing shift to be zero, whence
$v(z)=u(z)$. This last implication recovers contact for the original
supersolution.

\subsection{The exponential flux and weak inequalities}

On the full gradient space put
\[
 H(p,q)=q+\frac{|p|^2}{2},\qquad
 \mathcal A(p,q)=e^{H(p,q)/\tau}(p,1).
\]
Direct differentiation gives
\begin{equation}\label{eq:exp-flux}
 D\mathcal A(P)=\frac{e^{H(P)/\tau}}{\tau}A_\tau(p),
 \qquad
 \operatorname{div}\mathcal A(Df)
       =\frac{e^{H(Df)/\tau}}{\tau}G_\tau(Df,D^2f)
\end{equation}
for smooth $f$. This exponential structure is classical; see
\cite{SanchezMorgado}. It is not a new operator introduced here.
The derivative is symmetric positive definite. On $|P|\leq M$, one
may take
\[
 \begin{aligned}
 \lambda_0&=\bigl[\max\{2,(1+2M^2)/\tau\}\bigr]^{-1},
 &\Lambda_0&=M^2+1+\tau,\\
 \lambda_M&=e^{-M/\tau}\lambda_0/\tau,
 &\Lambda_M&=e^{(M+M^2/2)/\tau}\Lambda_0/\tau
 \end{aligned}
\]
in $\lambda_MI\leq D\mathcal A(P)\leq\Lambda_MI$.
Indeed,
$A_\tau(p)[(\zeta,s)]=(p\cdot\zeta+s)^2+\tau|\zeta|^2$;
the lower bound follows from
$s^2\leq2(s+p\cdot\zeta)^2+2M^2|\zeta|^2$.

\begin{lemma}\label{lem:flux-weak}
For a locally Lipschitz viscosity subsolution $f$,
$\operatorname{div}\mathcal A(Df)\geq0$ distributionally.
For a locally Lipschitz supersolution the inequality is reversed.
\end{lemma}
\begin{proof}
Fix a compact test support in a region where $f$ is $M$-Lipschitz,
with an additional compact margin. For small $\epsilon>0$ take the
full-variable sup convolution
\[
 f^\epsilon(z)=\sup_y\left(f(y)-\frac{|z-y|^2}{2\epsilon}\right).
\]
Optimizers satisfy $|y-z|\leq2M\epsilon$. On the smaller region
$f^\epsilon$ is $M$-Lipschitz, converges uniformly to $f$, and has
$D^2f^\epsilon\geq-I/\epsilon$ distributionally. Pulling an upper
test to an attaining fixed translate shows that $f^\epsilon$ is an
exact subsolution. This uses autonomy, not convexity in the gradient.
At its almost every Alexandrov twice-differentiability point,
\[
 \operatorname{tr}\bigl(D\mathcal A(Df^\epsilon)D^2f^\epsilon\bigr)
 \geq0.
\]
To see this, add a positive arbitrarily small quadratic to its
second-order Taylor polynomial and use the viscosity inequality.

Mollify $f^\epsilon$ on a smaller region, with mollifications
$f_j^\epsilon$. Then $|Df_j^\epsilon|\leq M$ and
$D^2f_j^\epsilon\geq-I/\epsilon$, giving the common lower bound
\[
 \operatorname{div}\mathcal A(Df_j^\epsilon)
 \geq-(m+1)\Lambda_M/\epsilon.
\]
Both gradients and Hessians converge almost everywhere to those of
$f^\epsilon$. For the Hessian statement, the distributional Hessian
of a semiconvex function is a matrix measure; its absolutely continuous
part is the Alexandrov Hessian, and mollifications of the singular part
converge to zero almost everywhere. This is the mollification fact in
\cite[Remark 3.3]{JulinJuutinen}, with the sign reversed from its
semiconcave formulation. For every nonnegative smooth compactly
supported $\phi$, integration by parts and one-sided Fatou give
\[
 \begin{split}
 -\int\mathcal A(Df^\epsilon)\cdot D\phi
 &=\lim_j\int\phi\,\operatorname{div}\mathcal A(Df_j^\epsilon)\\
 &\geq\int\phi\,
 \operatorname{tr}\bigl(D\mathcal A(Df^\epsilon)D^2f^\epsilon\bigr)
 \geq0.
 \end{split}
\]
In particular the displayed nonnegative density is locally integrable.
The favorable singular Hessian contribution has not been deleted.

It remains to pass to the nonlinear flux as $\epsilon\downarrow0$.
At a differentiability point $z$ of $f$, write an optimizer as
$y=z+\epsilon b_\epsilon$, where $|b_\epsilon|\leq2M$.
Taylor expansion, uniform for bounded $b_\epsilon$, and comparison
with $y=z+\epsilon Df(z)$ give
\[
 \tfrac12|b_\epsilon-Df(z)|^2=o(1).
\]
At a differentiability point of $f^\epsilon$, its gradient is
$b_\epsilon$. Along a fixed sequence $\epsilon_j\downarrow0$,
discard the null sets for $f$ and all its convolutions. Gradients then
converge almost everywhere and are uniformly bounded, so the smooth
flux converges in $L^1$ on the test support. This proves the first
assertion. For supersolutions use inf convolution, whose Hessian has
the upper bound $I/\epsilon$. Apply Fatou to the negative divergence
and compare its minimizing shift with $-\epsilon Df(z)$. This gives
the reversed inequality, without claiming that $-f$ satisfies the
same equation.
\end{proof}

\subsection{Contact propagation}

\begin{theorem}\label{thm:contact}
Let $\Omega\subset\mathbb R^{m+1}$ be connected and open. Suppose
$u\leq v$ are continuous, $u$ is a $G_\tau$-subsolution, $v$ is a
$G_\tau$-supersolution, and $u$ is locally uniformly semiconcave in
space. If they agree at one interior point, they agree everywhere.
\end{theorem}
\begin{proof}
First let both functions be locally Lipschitz and put $d=v-u\geq0$.
On a region where their gradients have norm at most $M$, define
\[
 B(z)=\int_0^1D\mathcal A\bigl(Du(z)+s(Dv(z)-Du(z))\bigr)\,ds.
\]
This is measurable, symmetric and uniformly elliptic between
$\lambda_MI$ and $\Lambda_MI$. Lemma~\ref{lem:flux-weak} and the
fundamental theorem of calculus imply
\[
 \mathcal A(Dv)-\mathcal A(Du)=B Dd,\qquad
 \int B Dd\cdot D\phi\geq0\quad(\phi\geq0).
\]
At a zero $z_0$ of $d$, choose any sufficiently small ball $B_r(z_0)$
compactly inside this region. Lax--Milgram gives $w\in H^1(B_r)$ with
$w-d\in H^1_0(B_r)$ and
$\int BDw\cdot D\phi=0$ for $\phi\in H^1_0(B_r)$.
The negative-part test gives $w\geq0$, because the boundary trace is
nonnegative. Testing the difference with $(w-d)^+$ gives
\[
 \int BD(w-d)^+\cdot D(w-d)^+\leq0,
\]
so $0\leq w\leq d$. Classical interior continuity and Harnack for
linear divergence-form weak solutions with bounded measurable
symmetric uniformly elliptic coefficients apply
\cite[Sections 2--3]{Crescimbeni}. For $m=1$ the total dimension is
two. If using a version stated only in dimensions above two,
extend $w$ constantly in a third coordinate and extend the coefficient
matrix to $\operatorname{diag}(B,1)$ on a product neighborhood.
Fubini verifies the weak equation, with common positive ellipticity;
the three-dimensional Harnack inequality then restricts to the
original plane. Thus $w(z_0)=0$ and Harnack,
followed by a chain of balls, gives $w=0$ throughout $B_r$.
Since $w-d\in H^1_0(B_r)$, $d$ has zero trace there. Its Lipschitz
continuity on a neighborhood of the closure identifies that trace
with its pointwise boundary restriction. Hence $d=0$ on the entire
sphere of radius $r$. The radius was arbitrary below a fixed margin,
so $d$ vanishes on a neighborhood of $z_0$.

For the stated pair, Lemma~\ref{lem:time-lipschitz} first makes $u$
locally Lipschitz. At every contact point apply the preceding cone
construction to $v$. The Lipschitz-pair argument makes $v_L=u$
nearby, and \eqref{eq:cone-lower} recovers $v=u$ there. The contact
set is consequently open. It is closed by continuity, nonempty by
hypothesis, and all of $\Omega$ by connectedness.
\end{proof}

\subsection{Geometric normalization and the affine-slice alternative}

We now prove Theorem~\ref{thm:viscous}. If $C_0=0$, fitted open boxes
are exactly affine and overlap to give one affine function. Suppose
$C_0>0$; it may be enlarged to at least one for estimates.
For a fitted frame $(x_0,t_0,R,p_0,a_0)$ the transformation
\begin{equation}\label{eq:app-frame}
 U(y,s)=\frac{V(x_0+R(y+p_0s),t_0+Rs)-V(x_0,t_0)}R
        -p_0\cdot y-\left(a_0+\frac{|p_0|^2}{2}\right)s
\end{equation}
preserves the fit constant, changing fit parameters to
$(p-p_0,a-a_0)$. Differentiation on test functions gives
$G_\tau[U]=R G_\tau[V]$, with $\tau$ unchanged. The equation-free
fit comparisons and compactness in \cite[Section 4]{OpenAI} therefore
apply: normalized functions have centered growth at most
$CK(1+\log K)^2$ on boxes of radius $K$, and form a locally compact
family. Locally uniform limits of these normalized frames retain
both $G_\tau=0$ and the all-center, all-radius fit
\eqref{eq:app-fit} with the same $C_0$. The equation passes by
viscosity stability on finite jets. For any fixed center and radius,
fit compactness gives a bounded sequence of fit coefficients;
extract a convergent subsequence and use strict fitted-domain margins
to pass the fit inequality, then continuity to include the boundary.
The auxiliary subsequence may depend on that center and radius,
since the limiting function is already fixed. This proves the
required fit for every center and radius without choosing one
subsequence for uncountably many fits.
The growth estimate supplies uniform spatial subquadratic tails on
compact time intervals.

If every horizontal slice is affine, write
$V(x,t)=P(t)\cdot x+c(t)$, with $P,c$ continuous. A fit at $(x,t)$
of radius $R$ gives $|p-P(t)|\leq C_0$ by taking $\ell=0$. For fixed
$\ell>0$ and $R\geq\ell$, evaluate the same fit at
$(h,\ell)=(p\ell,\ell)$ and $(-p\ell,-\ell)$ and add to obtain
\[
 \begin{split}
 \big|&[P(t+\ell)+P(t-\ell)-2P(t)]\cdot x\\
 &+\ell p\cdot[P(t+\ell)-P(t-\ell)]
      +c(t+\ell)+c(t-\ell)-2c(t)\big|\leq2C_0R.
 \end{split}
\]
All terms except the coefficient of $x$ are bounded uniformly as $x$
varies. Thus that coefficient vanishes, and continuity makes
$P(t)=bt+d$. The centered growth bound, applied at $x=\pm Ke$ and
$|t|\leq K$, yields $|P(t)|\leq C(1+\log K)^2$ and forces $b=0$.
Tests $P\cdot x+\psi(t)$ then reduce the equation for $c$ to
$\psi''=0$, so $c$ and $V$ are affine.

Otherwise a horizontal slice has a nonzero midpoint defect. The
transformation $V(x,t)\mapsto-V(x,-t)$ preserves the equation and
the fits, changes $(p,a)$ to $(-p,a)$, and reverses the defect. Choose
its sign positive.

\subsection{Extremal envelope and the changed equation}

Let $\kappa>0$ be the supremum, over all times and finite
positive-weight probability configurations, of
\[
 \frac{\sum_j\theta_jV(y_j,t)-V(\sum_j\theta_jy_j,t)}r,
 \qquad
 r^2=\sum_j\theta_j\left|y_j-\sum_k\theta_ky_k\right|^2>0.
\]
It is finite: normalize a radius-$r$ fit at the barycenter using
\eqref{eq:app-frame}. The centered growth bound is at most
$C(1+|y|^2)$; integration against a mean-zero, unit-second-moment
configuration bounds its normalized defect. Both the quotient and
$\tau$ are invariant under the fitted-frame transformation.

Take configurations tending to $\kappa$, normalized to mean zero and
second moment one, and locally uniformly convergent normalized
solutions $V_j\to\widehat V$. The corresponding measures have a weak
limit $\mu$ with mean zero and second moment at most one. First-moment
uniform integrability preserves the mean; uniform subquadratic tails
give uniform integrability for the varying value integrals. Hence
\[
 \int\widehat V(y,0)\,d\mu(y)-\widehat V(0,0)=\kappa.
\]
For $s>0$ use the barycentric envelope
\[
 \mathcal C_sV(x,t)=
 \sup_{\substack{\theta_j>0,\ \sum\theta_j=1\\
                   \sum\theta_jy_j=x}}
 \left\{\sum_j\theta_jV(y_j,t)
              -\frac1{2s}\sum_j\theta_j|y_j-x|^2\right\}.
\]
The Jensen bound, also valid in the normalized closure, gives
\begin{equation}\label{eq:envelope-sandwich}
 0\leq\mathcal C_sV-V\leq
 \sup_{r\geq0}\left(\kappa r-\frac{r^2}{2s}\right)
 =\frac{\kappa^2s}{2}.
\end{equation}
The measure $\mu$ is admissible for the envelope by its geometric
measure representation \cite[Section 5]{OpenAI}. At $s=1/\kappa$
its defect gives the reverse bound $\kappa/2$. Equality also forces
the second moment of $\mu$ to be one. Finite attainment, continuity
and the concave-envelope identity are the geometric conclusions of
that same lemma. A radius-$1/\kappa$ frame and time translation put
the equality at parameter one and time one. Rename the resulting
solution $V$, and set on $t>0$
\[
 W(x,t)=\mathcal C_tV(x,t)-\frac{\kappa^2t}{2}.
\]
Then
\begin{equation}\label{eq:extremal-contact}
 V-\frac{\kappa^2t}{2}\leq W\leq V,\qquad
 W(0,1)=V(0,1)=0,\qquad
 \mathcal C_1V(0,1)-V(0,1)=\frac{\kappa^2}{2}>0.
\end{equation}
The geometric identity
\[
 \mathcal C_tV(x,t)=\frac{|x|^2}{2t}
  +\operatorname{Conc}\left[V(\cdot,t)-\frac{|\cdot|^2}{2t}\right](x)
\]
gives $D^2_{xx}W\leq I/t$. Here $\operatorname{Conc}$ denotes the
least concave majorant. On each compact positive-time cylinder
$t\geq t_{\min}>0$, this gives spatial semiconcavity with the
single constant $1/t_{\min}$, uniformly in space and time. Thus,
once the subsolution property below is proved, the hypotheses of
Lemma~\ref{lem:time-lipschitz} and the contact theorem apply locally.

We must check the subsolution property for the changed PDE, rather
than import the zero-viscosity one. At an upper test $\phi$ of $W$,
choose a finite attaining configuration with positive weights and
barycenter $x$. The test for its weighted sum is
\[
 \Phi(Y,t)=\phi\left(\sum_j\theta_jy_j,t\right)
       +\frac{\kappa^2t}{2}
       +\frac1{2t}\sum_j\theta_j
                   \left|y_j-\sum_k\theta_ky_k\right|^2.
\]
Use the synchronized-time theorem of sums from
\cite[Sections 2 and 5]{OpenAI}. Divide each weighted closed jet by
its positive weight, apply its individual subsolution inequality,
and multiply the result back by that weight. Test the block matrix
inequality on the common-time vector $((q_j)_j,1)$, with
\[
 q_j=D_x\phi+\frac{y_j-x}{t},
\]
and on each of the $m$ common spatial-translation vectors
$((e_k)_j,0)$. First let the theorem-of-sums error tend to zero,
then let the time synchronization parameter tend to infinity. The
synchronization Hessian vanishes on all these vectors; large
individual time-gradient components do not enter $G_\tau$. The
resulting inequality is
\[
 D^2\Phi[((q_j)_j,1)]
       +\tau\sum_{k=1}^mD^2\Phi[((e_k)_j,0)]\geq0.
\]
Along the first vector, the centered offsets scale by $1+h/t$,
so their variance divided by $2(t+h)$ is affine in $h$. Along
common spatial translations the variance is constant. Their weighted
velocity is $D_x\phi$. Thus the displayed inequality is exactly
$G_\tau(D\phi,D^2\phi)\geq0$. It proves that $W$ is a subsolution,
including at zero horizontal gradient.

Theorem~\ref{thm:contact} now applies to \eqref{eq:extremal-contact}
on the connected positive-time half-space and gives $W=V$ there.
For any fixed $t>0$, subquadratic spatial growth implies that
$V(y,t)-|y|^2/(2t)$ attains a global maximum at some $y_*$. Average
its maximum inequality over any finite configuration with mean
$y_*$. The variance identity yields
\[
 \sum_j\theta_jV(y_j,t)
      -\frac1{2t}\sum_j\theta_j|y_j-y_*|^2\leq V(y_*,t).
\]
The singleton configuration attains equality, so
$\mathcal C_tV(y_*,t)=V(y_*,t)$. But $W=V$ requires the same gap to
be $\kappa^2t/2>0$ at every point, a contradiction. This proves
Theorem~\ref{thm:viscous}.
